\section{Music Theory Primer}\label{sec:background-music-theory}

To understand the abstractions provided by a musical domain-specific language, it is necessary to establish a
foundational understanding of western music theory.
Music is organized both vertically (pitch) and horizontally (time).

\subsection{Pitch, Notes, and Octaves}

The fundamental unit of melody and harmony is the \textit{note}, which represents a sound of a specific frequency.
In western music, the octave is divided into twelve semitones.
The seven natural notes are designated by the letters A through G.
\textit{Accidentals} modify these natural notes: a sharp (\#) raises the pitch by one semitone, while a flat
(b) lowers it by one semitone.

An \textit{octave} is the interval between one musical pitch and another with double its frequency.
To uniquely identify a note across the entire audible spectrum, it is denoted by its pitch class and its
octave number (e.g., \texttt{C4}, which is often referred to as "Middle C").
When two or more distinct notes are sounded simultaneously, they form a \textit{chord}, which provides the
harmonic foundation of a composition.

\subsection{Time, Rhythm, and Tempo}

While pitch defines the frequency of a sound, rhythm defines its placement and duration in time.
The speed at which a piece of music is played is governed by its \textit{tempo}, typically measured in Beats
Per Minute (BPM).
A tempo of 120 BPM signifies that there are 120 quarter-note beats in one minute, or one beat every 0.5 seconds.

Music is structurally divided into \textit{measures} (or bars) by a \textit{time signature}.
A time signature, such as 4/4, consists of two numbers: the numerator indicates the number of beats per
measure (four), and the denominator indicates the note value that constitutes one beat (the quarter note).
A \textit{rest} is a period of silence that has a specific duration, functioning as a "silent note" that
preserves the rhythmic alignment of the timeline.
